\section{Cap-Aware Width-Constrained PDB Selection}
\label{sec:supp-cap-selection}

This section specifies the focused implementation evaluated by the frozen
protocol.  It uses the safe finite-value ceiling from
Proposition~\ref{prop-transform}: for a PDB distance \(d<\infty\),

\[
  \phi_\kappa(d)=\min(d,\kappa),
  \qquad \phi_\kappa(\infty)=\infty.
\]

Thus each tested transform remains admissible and consistent, preserves
abstract dead ends, and cannot have larger cofactor width than its exact-value
parent.  The empirical comparison concerns whether selecting this transform
inside a fixed exact-width budget is useful; it does not modify the formal
safety results established earlier in this supplement.

\subsection{Shared Pattern Generators and Options}

Both compared selectors use the same abstract-state budget
\(B=100{,}000\), exact cofactor-width budget \(K=8\), fixed BDD variable
order, disabled dynamic reordering, CEGAR seed 2011, and
\texttt{cegar\_max\_time=10}.  Their fixed generator list, in protocol order,
is: the empty pattern; the longest BDD-order prefix fitting \(B\); the longest
goal/causal-order prefix fitting \(B\); a goal fill that skips a variable that
does not fit and continues down the same order; and one CEGAR pattern.  Patterns
are sorted and deduplicated while all generator-source labels are retained.
The empty pattern guarantees a feasible width-one member.

The two complete search strings are

\begin{sloppypar}
\small
\raggedright
\begin{description}
  \item[Exact \(K=8\)]
    \url{sym_fw_pdb(budget=100000,pattern_selection=exact_width_filter,cofactor_width_budget=8,cegar_max_time=10,cegar_seed=2011)};
  \item[Cap-aware \(K=8\)] the same string with
    \url{select_value_cap=true} appended.
\end{description}
\end{sloppypar}

There is one implementation-level normalization asymmetry that must not be
hidden by the shared generator description.  Only the cap-aware path deletes
domain-size-one variables from each generated pattern before sorting and pool
deduplication.  Such variables carry no state information, so this rewrite is
semantics preserving, but the exact-selector path is deliberately left
byte-for-byte unchanged.  Consequently, nominally identical generators and
options do \emph{not} by themselves establish identical raw candidate pools.
Any attribution to value capping is conditional on complete certified paired
traces having identical normalized uncapped-pool fingerprints in the stated
scope.

\subsection{Discrete Cap Ladder and Per-Pattern Choice}

The cap-aware selector materializes each raw PDB and its exact finite-value
histogram once.  It tests

\[
  \kappa\in\{0,1,2,4,8,16,32,64,128,256\}
\]

in increasing order, followed by the exact transform.  A finite cap at or
above that pattern's maximum finite distance is identical to exact values and
is omitted, so the emitted sequence is precisely the strict prefix below the
maximum followed by ``exact.''  All transforms reuse the one raw PDB;
infinity remains a separate dead-end set.

Feasibility means that the transformed total heuristic ADD has exact cofactor
width at most \(K\).  Lower caps are terminal relabelings of higher caps, so
width is monotone under this ordered refinement.  The retained representative
for a pattern is the \emph{last feasible tested member}: exact values when
feasible, otherwise the largest feasible ladder cap.  This strongest-feasible
step occurs before patterns are compared; the cap is not an extra cross-pattern
tie-breaker.

\subsection{Exact Width, Frozen Score, and Trace Certification}

For every tested member, the implementation constructs the same total ADD that
is reported for the selected heuristic.  Ordinary finite distances are
terminals, abstract dead ends receive a distinct terminal, and invalid binary
codes retain the default-zero extension.  Beginning with the root frontier, at
each manager level the exact-width procedure replaces residuals rooted at that
level by their two children, retains terminals and residuals that skip the
level, deduplicates by CUDD node identity, and records the largest frontier.
Primed transition variables do not occur in the heuristic ADD and therefore
only repeat a frontier.  Dynamic reordering is disabled, binding the measured
width to the order later used by search.

After taking one strongest feasible representative per pattern, the frozen
lexicographic score prefers: an initial-state dead-end proof; then larger
\(h(s_{\mathrm I})\); larger mean transformed finite distance; larger
abstract dead-end fraction; smaller measured cofactor width \(W\); fewer
abstract states; and the lexicographically smaller sorted pattern.  The means
and fractions are compared as exact integer ratios.  This is a deterministic
ranking surrogate, not a runtime-optimality claim.

\begin{proposition}[selector dominance]
\label{prop-selector-dominance}
Suppose the exact and cap-aware selectors complete on the same score-bearing
normalized raw-PDB pool with the same width budget. The cap-aware candidate set after
per-pattern selection contains every exact-feasible candidate. Hence its
winner has lexicographic score no smaller than the exact winner. It preserves
initial-state dead-end recognition and semantic nontriviality; if neither
winner recognizes the initial state as a dead end, its initial value is no
smaller. If the cap-aware winner is an exact endpoint, both selectors return
the same pattern and heuristic.
\end{proposition}

\begin{proof}
Consider any raw PDB whose exact transform has width at most $K$. Exact is
the final endpoint of that PDB's cap-aware ladder. Since it is feasible, the
last-feasible rule retains it as the per-pattern representative. It therefore
enters the cap-aware cross-pattern comparison with exactly the score it has in
the exact selector. This holds for every exact-feasible raw PDB, proving
candidate-set containment and the score inequality.

The first score component preserves initial-state dead-end recognition. If
the exact winner is semantic-nontrivial, it has either positive mean finite
distance or positive abstract dead-end fraction. A candidate ranked at least
as high either improves an earlier component---which already implies an
initial dead end or a positive finite value---or matches the earlier
components and cannot reduce the component witnessing nontriviality. Finally,
when both winners have a finite initial value, equality of the first component
forces the second component to be nondecreasing.

For the final statement, let $e$ and $p$ be the exact-selector and cap-aware
winners, and suppose $p$ is an exact endpoint. Candidate containment gives
$\mathrm{score}(p)\geq\mathrm{score}(e)$. Since $p$ also belongs to the exact
candidate set, optimality of $e$ gives the reverse inequality. The score is a
deterministic total order whose final pattern tie breaker is part of the
score-bearing pool identity, so $p=e$.
\end{proof}

This is a backward-compatibility result for the frozen information score, not
a coverage or runtime guarantee.

\begin{corollary}[selector safety]
\label{cor-selector-safety}
Under the assumptions of Theorem~\ref{thm-effort}, every completed selection
by either selector has effort at most $2K^2n$ times blind effort. If a selected
transform has finite cap $\kappa$, the sharper factor is
$2K\min\{K,\kappa+1\}n$. Thus $K=8$ gives at most $128n$ for every output,
and caps 2 and 4 give at most $48n$ and $80n$, respectively.
\end{corollary}

Every selected transform is admissible and consistent, passes the measured
width test $W\leq K$, and has at most $W\leq K$ finite values by
Proposition~\ref{prop-range}. A cap has at most $\kappa+1$ finite values.
Substitution in Corollary~\ref{cor-effort} proves the claim. This certificate
does not require paired-pool identity and does not bound runtime or relational
products.

The cap-aware whole-trace certificate records each tested transform, its cap,
feasibility decision, transformed \(W\), ADD inner nodes \(A\), terminal count
\(T\), total ADD nodes \(U=A+T\), finite-value count \(V\), and the raw
counterparts.  The exact raw histogram is emitted once per materialized
pattern.  The fail-closed parser reconstructs the capped initial value, finite
sum, and value count; checks the cap order, within-pattern raw identity,
monotonicity of \(W,A,T,U,V\), the last-feasible representative, the frozen
score winner, and agreement among the selected record, final heuristic record,
and whole-trace hashes.  Cross-protocol raw-pool equality is then a separate
paired certification condition, for the domain-one normalization reason above.

\section{Focused Development Gate and Full Evaluation Protocol}
\label{sec:supp-cap-protocol}

\subsection{Outcome-Informed Development Screen}

The launch gate used a fixed 50-task, 25-domain development manifest with two
tasks per domain.  Its composition had been assembled during earlier
development using archived outcome summaries to enrich difficult or
discriminating cases.  It is therefore neither a held-out sample nor an
independent population estimate.  None of those archived measurements is a
cell in this study.

All six gate configurations were rerun under one pinned planner, protocol, and
environment: blind forward search; the \(B=100{,}000\) CEGAR PDB; exact and
cap-aware selectors at \(K=8\); and exact and cap-aware selectors at \(K=32\).
This is a new \(50\times6=300\)-cell mechanism screen.  \(K=8\) is the fixed
primary comparison; \(K=32\) is a non-promotable sensitivity and cannot replace
it.

The Boolean promotion rule opens the full launch gate if and only if all of the
following hold:

\begin{enumerate}
  \item all 300 screen cells are present and validate;
  \item all 50 exact-versus-cap-aware \(K=8\) pairs have complete certified
    traces and identical normalized raw-pool fingerprints;
  \item every solved cost agrees across screen configurations, with no
    solved/proved-unsolvable contradiction;
  \item the cap-aware \(K=8\) selector has strictly more semantically
    nontrivial selections than exact \(K=8\), where nontrivial means a positive
    finite-distance sum or a positive abstract-dead-end count; and
  \item cap-aware \(K=8\) has coverage wins minus losses at least zero.
\end{enumerate}

The gate is conjunctive, contains no tunable threshold beyond those stated,
and never consults \(K=32\) for promotion.  The full protocol recomputes the
digest-pinned gate certificate before launch rather than accepting a manually
transcribed decision.

\subsection{Frozen Population, Matrix, and Environment}

The source inventory contains 1697 optimal sequential IPC tasks through 2023.
Disjoint frozen exclusions remove 290 tasks with zero-cost operators and 30
with normalized axioms, leaving 1377 positive-cost, axiom-free tasks in 46
domains.  The exclusions align the empirical population with the positive-cost
assumption of Lemma~\ref{lem-identity} and the abstraction implementation's
axiom restriction.  Task order, source bytes, operator-cost support, normalized
axiom status, and the complete source manifest are digest-attested before
launch.

The full matrix is exactly \(1377\times5=6885\) cells:

\begin{enumerate}
  \item blind forward symbolic search;
  \item causal-graph linear merge-and-shrink with
    \url{max_states=10000,value_cap=-1,align_merge_order=false};
  \item the \(B=100{,}000\), seed-2011, 10-second-option CEGAR PDB;
  \item the exact-value \(K=8\) selector; and
  \item the cap-aware \(K=8\) selector.
\end{enumerate}

The first three are descriptive context; the last two define the primary
contrast.  All five use the same pinned \texttt{release\_no\_lp} binary and
preprocessor, the same task-specific fixed BDD order, one CPU on homogeneous
AMD EPYC 9755 nodes, an aggregate-process CPU limit of 300 seconds, and an
8-GiB planner memory limit.  Construction is inside those limits and is
charged to planner time.  There is one run per cell.  The scheduler groups at
most seven runs per array task inside a reviewed 45-minute envelope with 9 GiB
per CPU; this envelope is not represented as a per-cell wall-clock watchdog.

Every retained cell must have one atomic completion marker and a recognized,
reconciled terminal outcome.  Missing, duplicate, starved, unrecognized, or
infrastructure-failed cells invalidate the census; there is no post-hoc cell
fill.  Solved costs must agree across all five configurations, and no solved
record may conflict with a proved-unsolvable record.  Micro-PAR2 is the
arithmetic mean of planner CPU seconds with 600 seconds assigned to every
unsolved cell.

\subsection{Primary Estimand and Prespecified Sensitivity}

The primary population is the frozen 1327-task complement obtained by deleting
the 50 development tasks from the full census while retaining all 46 domains.
The primary estimand is cap-aware minus exact \(K=8\) equal-domain macro
coverage: compute each selector's coverage within each domain of this
complement, average the 46 domain rates equally, and subtract exact from
cap-aware.  The paired task-micro disclosure reports cap-aware wins, losses,
both solved, both unsolved, and discordance on the same fixed complement.

The sole prespecified sensitivity repeats the contrast on all 1377 tasks,
including the development set.  It is a fixed-population census sensitivity,
not a replacement primary analysis.  Neither scope uses confidence intervals,
\(p\)-values, or population-generalization claims.  Mechanism attribution also
requires complete certified traces and identical raw-pool fingerprints for
every compared task in the scope; selection differences are required before a
mechanism claim.

\subsection{Secondary Descriptive Contract and Its Timing}

The primary estimand and launch rule were frozen before the development
aggregate was inspected.  The secondary descriptive contract was frozen later,
while full execution was active but before full properties were fetched and
before any scientific outcome or aggregate was read; only completion-marker
and error-file health monitoring had begun.  Its version-2 amendment added
scope-specific solved counts and solved-runtime totals solely so every reported
micro-PAR2 value can be reconstructed.  The contract is explicitly
secondary-descriptive and cannot change promotion, the primary scope, or the
primary estimand.

For both the 1327-task primary complement and the 1377-task sensitivity, the
contract reports complete-certified selector-trace counts; semantic
nontriviality; effective-cap and first-provenance-source histograms; and
pattern-size, \(W,A,T,U,V\) summaries.  Paired operational descriptions use
cap-aware over exact and state their own eligible denominator: fixed-scope
micro-PAR2; jointly solved positive-time CPU geometric mean; ratios of image
time and expanded-BDD-node totals on the same complete-certified image pairs;
and construction-time totals on pairs with observed completed constructions.
The five-configuration full-census context additionally reports solved counts,
micro-PAR2, solved-runtime summaries, and completed-construction summaries.
Missing observations are never imputed.

\section{Focused Full-Evaluation Results}
\label{sec:supp-cap-tables}

\CapPrimaryText\ \CapSensitivityText\ Under the prespecified all-scope rule,
\CapMechanismText\ The missing 21 pairs are not instrumentation losses: their
runs ended before both selectors produced completed traces under the stated
time or memory limits. On all \CapPrimaryTracePairs\ completed pairs, the raw
pools match, so Proposition~\ref{prop-selector-dominance} applies. Cap-aware
selection adds \CapPrimarySemanticWins\ semantic-nontrivial winners and loses
\CapPrimarySemanticLosses. This is a complete conditional result on the
reported paired denominator, while the frozen global mechanism label remains
not certified. \CapScopeCaveat

\begin{table}[htbp]
  \centering
  \small
  \caption{Prespecified all-task census sensitivity. The legacy study's
    primary contrast excludes the 50 development tasks.}
  \label{tab:supp-cap-census}
  \begin{tabular}{@{}lrr@{}}
    \toprule
    Selector & Equal-domain macro coverage (\%) & Solved/tasks \\
    \midrule
    \CapCensusContrastRows
    \bottomrule
  \end{tabular}
\end{table}

\CapPrimarySecondaryText\ \CapRatioCaveat\ \CapCensoringCaveat

\begin{table}[htbp]
  \centering
  \footnotesize
  \setlength{\tabcolsep}{3.5pt}
  \caption{Paired operational summaries for the primary complement and the
    all-task sensitivity. The final column is cap-aware minus exact for
    micro-PAR2 and cap-aware over exact for every ratio; dashes mark cells that
    are not separately meaningful for a geometric-mean ratio.}
  \label{tab:supp-cap-operational}
  \begin{tabular}{@{}lrrrr@{}}
    \toprule
    Metric & Eligible & Cap-aware & Exact & $\Delta$/ratio \\
    \midrule
    \multicolumn{5}{@{}l}{\textit{Primary development complement}} \\
    \CapPrimarySecondaryRows
    \midrule
    \multicolumn{5}{@{}l}{\textit{All-task census sensitivity}} \\
    \CapCensusSecondaryRows
    \bottomrule
  \end{tabular}
\end{table}

\begin{table}[htbp]
  \centering
  \small
  \caption{Completed-construction observations for the five full-census
    configurations. Missing constructions are not imputed.}
  \label{tab:supp-cap-construction}
  \begin{tabular}{@{}lrr@{}}
    \toprule
    Configuration & Observed/tasks & Total seconds \\
    \midrule
    \CapFullConstructionRows
    \bottomrule
  \end{tabular}
\end{table}

\begin{table}[p]
  \centering
  \small
  \setlength{\tabcolsep}{3.5pt}
  \caption{Selector representation summaries. \(A\) is ADD inner nodes,
    \(T\) terminals, \(U=A+T\), and \(V\) finite values. Each metric carries
    its observed-trace denominator.}
  \label{tab:supp-cap-selector-summary}
  \begin{tabular}{@{}lllrrrrr@{}}
    \toprule
    Scope & Selector & Metric & Observed/traces & Min & Median & Max & Total \\
    \midrule
    \CapSelectorSummaryRows
    \bottomrule
  \end{tabular}
\end{table}

\begin{landscape}
\begin{table}[p]
  \centering
  \footnotesize
  \setlength{\tabcolsep}{3pt}
  \caption{First-provenance source and effective-cap histograms. Counts are
    shown against the complete-certified trace denominator of that scope and
    selector.}
  \label{tab:supp-cap-histograms}
  \begin{minipage}[t]{0.48\linewidth}
    \centering
    \textbf{First-provenance source}\\[2pt]
    \begin{tabular}{@{}lllr@{}}
      \toprule
      Scope & Selector & Source & Count/traces \\
      \midrule
      \CapSelectorSourceRows
      \bottomrule
    \end{tabular}
  \end{minipage}\hfill
  \begin{minipage}[t]{0.48\linewidth}
    \centering
    \textbf{Effective cap}\\[2pt]
    \begin{tabular}{@{}lllr@{}}
      \toprule
      Scope & Selector & Cap & Count/traces \\
      \midrule
      \CapEffectiveCapRows
      \bottomrule
    \end{tabular}
  \end{minipage}
\end{table}
\end{landscape}

\subsection{Post-Hoc Descriptive Diagnostics}

After completing the prospective analysis, we added a separate
digest-pinned audit to answer the adversarial questions of whether the
aggregate operational reduction is specific to actual cap interventions and
whether it is dominated by a single domain. Its role is
\CapPosthocAnalysisRole, with inferential policy
\CapPosthocInferencePolicy: it changes neither the frozen estimand nor the
secondary contract and reports no tests or intervals.

The trace strata are exhaustive on each complete-certified paired denominator.
All \CapPosthocPrimaryExactEndpointPairs\ primary exact endpoints have zero
selected-heuristic identity mismatches. All
\CapPosthocPrimaryFinitePairs\ primary finite-cap endpoints materially change
the chosen raw PDB and have raw width above $K=8$; their raw-width range is
\CapPosthocPrimaryRawWidthMinimum--\CapPosthocPrimaryRawWidthMaximum, with
median \CapPosthocPrimaryRawWidthMedian. Among
\CapPosthocPrimaryInitialNondeadPairs\ primary pairs without an initial
dead-end proof, the cap-aware initial value is higher/equal/lower on
\CapPosthocPrimaryInitialHigher/\CapPosthocPrimaryInitialEqual/
\CapPosthocPrimaryInitialLower\ tasks.

The paired coefficient audit avoids cross-pattern confounding by evaluating
the chosen cap-aware PDB immediately before and after its terminal transform.
For every primary finite-cap trace, the coarse certificate coefficient
$c=2\values W$ decreases strictly. Table~\ref{tab:supp-posthoc-compression}
reports both the marginal coefficients and the exact paired raw-to-capped
ratios; these are structural certificates, not observed overhead ratios.

\begin{table}[htbp]
  \centering
  \caption{Post-hoc same-PDB certificate compression on every finite-cap
    trace. P95 uses nearest rank. ``Strict'' counts strict paired coefficient
    reductions.}
  \label{tab:supp-posthoc-compression}
  \resizebox{\textwidth}{!}{%
  \begin{tabular}{@{}lrrrrrrrrrrr@{}}
    \toprule
    Scope & Pairs & \multicolumn{3}{c}{Raw $c$} &
      \multicolumn{3}{c}{Capped $c$} &
      \multicolumn{3}{c}{Raw/capped} & Strict \\
    & & Median & P95 & Max & Median & P95 & Max & Median & P95 & Max & \\
    \midrule
    \CapPosthocCoefficientCompressionRows
    \bottomrule
  \end{tabular}
  }
\end{table}

\begin{table}[htbp]
  \centering
  \small
  \caption{Post-hoc exhaustive strata of complete-certified paired selector
    traces. Coverage columns retain every paired trace, including timeouts.}
  \label{tab:supp-posthoc-trace-strata}
  \begin{tabular}{@{}llrrrrr@{}}
    \toprule
    Scope & Endpoint & Pairs & Wins & Losses & Both solved & Both unsolved \\
    \midrule
    \CapPosthocTraceStrataRows
    \bottomrule
  \end{tabular}
\end{table}

\begin{table}[htbp]
  \centering
  \caption{Post-hoc operational localization. Ratios are cap-aware over exact;
    L/E/H counts pairs on which the cap-aware metric is lower/equal/higher.
    Eligibility requires complete certified schema-v2 image metrics, so these
    rows cannot explain coverage discordances or joint timeouts.}
  \label{tab:supp-posthoc-operational}
  \resizebox{\textwidth}{!}{%
  \begin{tabular}{@{}lllrrrrrrr@{}}
    \toprule
    Scope & Endpoint & Metric & Eligible & Cap-aware total & Exact total & Ratio
      & L & E & H \\
    \midrule
    \CapPosthocOperationalRows
    \bottomrule
  \end{tabular}
  }
\end{table}

On the primary operational subset, every eligible exact-endpoint node count is
identical, whereas finite caps give the full node reduction. Image
time is a measured duration and therefore need not be identical across reruns
of the same operations. The exact-endpoint coverage losses in
Table~\ref{tab:supp-posthoc-trace-strata} likewise occur at the resource limit
despite identity of the selected heuristic, so they are not evidence against
Proposition~\ref{prop-selector-dominance}.

\begin{table}[htbp]
  \centering
  \caption{Post-hoc conservation across completed buckets. ``Pos./zero''
    splits pairs by whether the bucket denominator is positive or zero in both
    configurations. Bucket, node/bucket and image/bucket ratios are cap-aware
    over exact; the last columns give taskwise L/E/H signs for the normalized
    metrics on positive-denominator pairs.}
  \label{tab:supp-posthoc-per-bucket}
  \resizebox{\textwidth}{!}{%
  \begin{tabular}{@{}llrrrrrrrrll@{}}
    \toprule
    Scope & Endpoint & Pairs & Pos. & Zero & Cap buckets & Exact buckets &
      Buckets & Nodes/bucket & Image/bucket & Node L/E/H & Image L/E/H \\
    \midrule
    \CapPosthocOperationalPerBucketRows
    \bottomrule
  \end{tabular}
  }
\end{table}

These operational comparisons include pattern changes induced by selection;
the same-PDB coefficient audit above is the isolated terminal-transform
comparison. Finite caps can refine the search into more completed buckets while
reducing the aggregate representation and image work per bucket. Exact
endpoints are the negative control: their completed-bucket totals, node totals
and node-per-bucket values are identical. Image time need not be identical
because it includes measured systems variation around the same operations.

The target-metric audit uses the exact expansion-size measure bounded by
Theorem~\ref{thm-effort}. Eligibility requires complete identical-pool selector traces,
certified solved searches for blind, exact and cap-aware configurations, one
BDD piece per expanded bucket, agreement of all three solution costs,
schema-v2 metrics, and
integral effort equal to the validated expanded-BDD-node counter. In
particular, zero/zero unsolvable operational pairs do not enter this
finite-$\copt$ comparison.

\begin{table}[htbp]
  \centering
  \caption{Post-hoc target-metric effort on triple-solved tasks. Ratios use
    pooled totals; L/E/H compares cap-aware with exact taskwise.}
  \label{tab:supp-posthoc-effort}
  \resizebox{\textwidth}{!}{%
  \begin{tabular}{@{}llrrrrrrrl@{}}
    \toprule
    Scope & Endpoint & Tasks & Blind & Exact & Cap-aware & Cap/blind &
      Exact/blind & Cap/exact & Cap/exact L/E/H \\
    \midrule
    \CapPosthocTripleSolvedEffortRows
    \bottomrule
  \end{tabular}
  }
\end{table}

\begin{table}[htbp]
  \centering
  \small
  \caption{Post-hoc leave-one-domain-out stability of the finite-cap
    triple-solved effort ratios. Parentheses name the omitted domain at each
    extreme.}
  \label{tab:supp-posthoc-effort-lodo}
  \begin{tabular}{@{}llrrr@{}}
    \toprule
    Scope & Ratio & Base & Minimum & Maximum \\
    \midrule
    \CapPosthocFiniteEffortLodoRows
    \bottomrule
  \end{tabular}
\end{table}

The effort comparison evaluates the theorem's target measure, not its
profile-sensitive numerical upper bound: the logs do not retain both cofactor
profiles at every layer. Its triple-solved denominator also cannot explain
coverage discordances or joint timeouts.

\begin{table}[htbp]
  \centering
  \small
  \caption{Post-hoc leave-one-domain-out descriptions. ``Base'' is the
    unmodified scope; the last columns show the extremal omission and the
    omitted domain. Coverage is in percentage points and nodes are ratios.}
  \label{tab:supp-posthoc-lodo}
  \begin{tabular}{@{}llrrr@{}}
    \toprule
    Scope & Metric & Base & Minimum & Maximum \\
    \midrule
    \CapPosthocLeaveOneDomainOutRows
    \bottomrule
  \end{tabular}
\end{table}

\begin{table}[htbp]
  \centering
  \small
  \caption{Post-hoc distribution of the selected universal coefficient
    $c=2VW$ in the certificate $c n$, computed per completed selector trace.
    This is a conservative structural certificate, not measured search
    overhead.}
  \label{tab:supp-posthoc-certificate}
  \begin{tabular}{@{}llrrrrr@{}}
    \toprule
    Scope & Selector & Traces & Min & Median & P95 & Max \\
    \midrule
    \CapPosthocCertificateCoefficientRows
    \bottomrule
  \end{tabular}
\end{table}

\begin{table}[htbp]
  \centering
  \small
  \caption{Post-hoc distribution of the sharper finite-cap coefficient
    $2K\min\{K,\kappa+1\}$. Selected caps yielding the same coefficient are
    grouped. These counts were not used to choose the cap ladder.}
  \label{tab:supp-posthoc-cap-certificate}
  \begin{tabular}{@{}lrrl@{}}
    \toprule
    Scope & Coefficient & Selections & Selected caps \\
    \midrule
    \CapPosthocUniversalCoefficientRows
    \bottomrule
  \end{tabular}
\end{table}
